\begin{frame}{Ошибка отношения аппроксимации к КФ}
    При фиксированной кривой $A=F(q)$ и измерении $x=C(q)$:
    \[
        f=A/x,\qquad \sigma_f=\left|\frac{\partial f}{\partial x}\right|\sigma_x
        =\frac{|A|}{x^2}\sigma_x.
    \]
    Например, $A=1.5$, $x=1.2$, $\sigma_x=0.06$:
    $f=1.25$, $\sigma_f=0.0625$.
    \begin{itemize}
        \item Ошибка вычисленной кривой явно равна нулю;
              ошибки отношения определяются измеренной КФ.
        \item Это условная диагностика при фиксированных параметрах.
              Неопределённость параметров фита и их связь с данными
              требуют отдельного расчёта ковариации.
        \item При $x=0$ отношение не определено; такой бин не интерпретируется.
    \end{itemize}
\end{frame}

\begin{frame}{Численное влияние округления суммы весов}
    \small
    Контрольный пример: $n=10000$, $A=10000.0049$,
    $Q_A=A^2/n+0.01$. Округляется только $A$, $Q_A$ сохраняется.
    \[
        \sigma_C=\sqrt{\frac{Q_A-A^2/n}{n(n-1)}}.
    \]
    \begin{center}
        \begin{tabular}{c|c|c|c}
            Знаков после точки & $C=A/n$ & $Q_A-A^2/n$ & $\sigma_C$\\ \hline
            2 & 1.00000000 & 0.0198000 & $1.40720\cdot10^{-5}$\\
            3 & 1.00000050 & 0.0098000 & $0.99000\cdot10^{-5}$\\
            4 & 1.00000049 & 0.0100000 & $1.00005\cdot10^{-5}$\\
            6 & 1.00000049 & 0.0100000 & $1.00005\cdot10^{-5}$
        \end{tabular}
    \end{center}
    При двух знаках $C$ меняется всего на $4.9\cdot10^{-7}$,
    а ошибка увеличивается на $40.7\%$.
    Округлять для вывода следует после расчёта ошибок и $\chi^2$.
    Это эксперимент с десятичным округлением, не модель накопления \texttt{TH3F}.
\end{frame}

\begin{frame}{Накопление одинаковых весов в ROOT: float и double}
    \small
    Один бин, 10000 заполнений; веса чередуются:
    $1.00100049$ и $0.99900049$. \texttt{Sumw2} включён до заполнения.
    \begin{center}
        \begin{tabular}{l|c|c|c}
            Тип & $A$ & $Q_A-A^2/n$ & Сырая $\sigma_C$\\ \hline
            \texttt{TH1F} & 10000 & 0.0198000 & $1.40720\cdot10^{-5}$\\
            \texttt{TH1D} & 10000.0049 & 0.0100000 & $1.00005\cdot10^{-5}$
        \end{tabular}
    \end{center}
    \begin{itemize}
        \item Данные и порядок заполнения одинаковы; различается тип хранения суммы.
        \item \texttt{Sumw2} хранится в double в обоих случаях.
        \item Аналогичный механизм действует в ячейках \texttt{TH3F}/\texttt{TH3D}.
        \item В программе неразрешимая относительно границы округления
              дисперсия исключается из фита. Таблица показывает сырую формулу
              до этой проверки, а не используемую программой ошибку.
    \end{itemize}
    Воспроизведение: \texttt{root -l -b -q tools/statistics\_review.C}.
\end{frame}

\begin{frame}{Ручной расчёт $\chi^2/\mathrm{ndf}$}
    \small
    Контрольные измерения имеют $\sigma_i=0.1$.
    Подгоняется одна константа $F=\overline C=1.2$.
    \begin{center}
        \begin{tabular}{c|ccccc}
            $C_i$ & 1.0 & 1.1 & 1.2 & 1.3 & 1.4\\ \hline
            $(C_i-F)/\sigma_i$ & $-2$ & $-1$ & 0 & 1 & 2\\
            Вклад в $\chi^2$ & 4 & 1 & 0 & 1 & 4
        \end{tabular}
    \end{center}
    \[
        \chi^2=4+1+0+1+4=10,\quad
        \mathrm{ndf}=5-1=4,\quad \chi^2/\mathrm{ndf}=2.5.
    \]
    В 3D-анализе суммируются допустимые ячейки области фита;
    $\mathrm{ndf}=N_{\mathrm{used}}-N_{\mathrm{free}}$.
    Замороженные параметры не вычитаются. Для интегрального фита
    $F_i$ является средним модели внутри ячейки.
    Пример проверяет алгоритм; это не результат для данных столкновений.
\end{frame}
